\documentclass[12pt]{article}
\usepackage[utf8]{inputenc}
\usepackage[T1]{fontenc}
\usepackage{hyperref}
\title{Pintiini Language Resources for K-12 Curriculum Development}
\author{Compiled by Claude Sonnet 5.5 (claude-sonnet-5-5) for Andr\'as Kornai}
\date{}
\begin{document}
\maketitle

\section{Introduction}
Pintiini (ISO 639-3 \texttt{pti}; Glottocode pint1251; also spelled Pindiini) is a Western Desert (Pama-Nyungan) variety of Western Australia's Goldfields region. AIATSIS notes that Pindiini, Wangkatja and Wangkatha are very close and hard to distinguish, so the resources below are for the closely related Wangkatja/Wangkatha rather than Pintiini alone. I found no Pintiini-specific school materials.

\section{Primary Educational Resources}
\begin{itemize}
\item Goldfields Aboriginal Language Centre (Wangka) 2025 Wangkatja grammar teaching handouts: Questions in Wangkatja, Pronouns, and Noun Suffixes (V1.2), and Verb Suffixes. \url{https://wangka.com.au/wp-content/uploads/2025/05/Questions-in-Wangkatja-V1.2-.pdf} \\ \url{https://wangka.com.au/wp-content/uploads/2025/05/Wangkatja-Pronouns-V1.2.pdf} \\ \url{https://wangka.com.au/wp-content/uploads/2025/05/Wangkatja-Noun-Suffixes-V1.2.pdf} \\ \url{https://wangka.com.au/wp-content/uploads/2025/03/Wangkatja-Verb-Suffixes-I-IMP-and-Tense-suffixes.pdf}
\end{itemize}

\section{Linguistic and Reference Materials}
\begin{itemize}
\item AUSTLANG entry A103 (Wangkatja), AIATSIS. \url{https://collection.aiatsis.gov.au/austlang/language/A103}
\item AUSTLANG entry A12, AIATSIS, which the search summary tied to the Wangkatha name and the 2002 \emph{Wangkatha dictionary} (Wangkanyi Ngurra Tjurta, Kalgoorlie); I did not see the dictionary itself online. \url{https://collection.aiatsis.gov.au/austlang/language/A12}
\end{itemize}

\section{Religious and Community Resources}
\begin{itemize}
\item Goldfields Aboriginal Language Centre (Wangka) resource catalogue. \url{https://wangka.com.au/product-category/resources/page/3}
\end{itemize}

\section{Digital and Audio Resources}
No Pintiini-specific audio or video was found in search results.

\section{Curriculum Development Context}
The language is not known to be taught in schools. A curriculum would need to start with the Goldfields language centre and local community, building on Wangkatja materials and related Ngaanyatjarra resources, and settling with elders what counts as Pintiini versus Wangkatja.

\end{document}